\subsection{THE ORDERED GROUP OF RATIONAL NUMBERS}

We have defined the ordering \(x \leqslant y\) on the set Q of rational numbers; we have seen that this ordering makes Q a linearly ordered set, and that it is compatible with the additive group structure of Q, i.e.~for each \(z \in Q\) the relation \(x \leqslant y\) is equivalent to \(x + z \leqslant y + z\) (that is, the ordering is invariant under translations). We recall the notation (which is used in any linearly ordered group)

\[
\begin{array}{l} x ^ {+} = \sup (x, 0), \\ x ^ {-} = \sup (- x, 0) = (- x) ^ {+}, \\ | x | = \sup (x, - x); \end{array}
\]

\(|x|\) is called the absolute value of \(x\) , and we have

\[
x = x ^ {+} - x ^ {-}, \quad | x | = x ^ {+} + x ^ {-}
\]

and the triangle inequality

\[
| x + y | \leqslant | x | + | y |,\tag{1}
\]

together with the inequality

\[
\mid | x \mid - \mid y \mid \mid \leqslant \mid x - y \mid\tag{2}
\]

which is an immediate consequence of (1); moreover

\[
\left| x ^ {+} - y ^ {+} \right| \leqslant | x - y |.\tag{3}
\]

The relations \(x \geqslant 0\) , \(x = x^{+}\) , \(x^{-} = 0\) , \(|x| = x\) (resp. \(x \leqslant 0\) , \(x = -x^{-}\) , \(x^{+} = 0\) , \(|x| = -x\) ) are equivalent. The relation \(|x| = 0\) is equivalent to \(x = 0\) ; if \(a \geqslant 0\) , the relation \(|x| \leqslant a\) is equivalent to \(-a \leqslant x \leqslant a\) , and the relation \(|x| \geqslant a\) is equivalent to `` \(x \geqslant a\) or \(x \leqslant -a\) ''. For all \(x, y\) , in \(\mathbf{Q}\) , we have

\begin{enumerate}
\def\labelenumi{(\arabic{enumi})}
\setcounter{enumi}{3}
\tightlist
\item
\end{enumerate}

\[
\sup (x, y) + z = \sup (x + z, y + z),\tag{4}
\]

\[
\inf (x, y) = - \sup (- x, - y),\tag{5}
\]

and, as particular cases,

\begin{enumerate}
\def\labelenumi{(\arabic{enumi})}
\setcounter{enumi}{5}
\tightlist
\item
\end{enumerate}

\[
\sup (x, y) = x + (y - x) ^ {+} = x + (x - y) ^ {-},\tag{6}
\]

\[
\inf (x, y) = x - (y - x) ^ {-} = x - (x - y) ^ {+},\tag{7}
\]

Finally, let \(Q_{+}\) denote the set of rational numbers \(\geqslant0\) ; we have then

\begin{enumerate}
\def\labelenumi{(\arabic{enumi})}
\setcounter{enumi}{7}
\tightlist
\item
\end{enumerate}

\[
\mathbf {Q} _ {+} + \mathbf {Q} _ {+} \subset \mathbf {Q} _ {+},\tag{8}
\]

\[
\mathbf {Q} _ {+} \cap (- \mathbf {Q} _ {+}) = \{\mathrm{o} \},\tag{9}
\]

\[
\mathbf {Q} _ {+} \cup (- \mathbf {Q} _ {+}) = \mathbf {Q},\tag{10}
\]

The relation \(x \leqslant y\) is equivalent to \(y - x \in \mathbf{Q}_{+}\) .

We shall use this ordering to define a topology on Q compatible with its additive group structure.

